\documentclass[10pt,a4paper]{amsart}
\usepackage{amsmath,amssymb,amsthm,graphicx}
% For dealing with references we use the comment environment
\usepackage{verbatim}
\newenvironment{reference}{\comment}{\endcomment}
%\newenvironment{reference}{}{}
\newenvironment{slogan}{\comment}{\endcomment}
\newenvironment{history}{\comment}{\endcomment}

% For commutative diagrams we use Xy-pic
\usepackage[all]{xy}
% Theorem environments.
%
\theoremstyle{plain}
\newtheorem{theorem}[subsection]{Theorem}
\newtheorem{proposition}[subsection]{Proposition}
\newtheorem{lemma}[subsection]{Lemma}

\theoremstyle{definition}
\newtheorem{definition}[subsection]{Definition}
\newtheorem{example}[subsection]{Example}
\newtheorem{exercise}[subsection]{Exercise}
\newtheorem{situation}[subsection]{Situation}

\theoremstyle{remark}
\newtheorem{remark}[subsection]{Remark}
\newtheorem{remarks}[subsection]{Remarks}

\numberwithin{equation}{subsection}

% Macros
%
\def\lim{\mathop{\mathrm{lim}}\nolimits}
\def\colim{\mathop{\mathrm{colim}}\nolimits}
\def\Spec{\mathop{\mathrm{Spec}}}
\def\Hom{\mathop{\mathrm{Hom}}\nolimits}
\def\Ext{\mathop{\mathrm{Ext}}\nolimits}
\def\SheafHom{\mathop{\mathcal{H}\!\mathit{om}}\nolimits}
\def\SheafExt{\mathop{\mathcal{E}\!\mathit{xt}}\nolimits}
\def\Sch{\mathit{Sch}}
\def\Mor{\mathop{\mathrm{Mor}}\nolimits}
\def\Ob{\mathop{\mathrm{Ob}}\nolimits}
\def\Sh{\mathop{\mathit{Sh}}\nolimits}
\def\NL{\mathop{N\!L}\nolimits}
\def\CH{\mathop{\mathrm{CH}}\nolimits}
\def\proetale{{pro\text{-}\acute{e}tale}}
\def\etale{{\acute{e}tale}}
\def\QCoh{\mathit{QCoh}}
\def\Ker{\mathop{\mathrm{Ker}}}
\def\Im{\mathop{\mathrm{Im}}}
\def\Coker{\mathop{\mathrm{Coker}}}
\def\Coim{\mathop{\mathrm{Coim}}}

% Boxtimes
%
\DeclareMathSymbol{\boxtimes}{\mathbin}{AMSa}{"02}

%
% Macros for moduli stacks/spaces
%
\def\QCohstack{\mathcal{QC}\!\mathit{oh}}
\def\Cohstack{\mathcal{C}\!\mathit{oh}}
\usepackage{cleveref}
\crefname{originalsection}{original section}{original sections}
\Crefname{originalsection}{Original section}{Original sections}
\makeatletter
\newlabel{categories-section-localization}{{04VB}{source}{}{sourcepointer-1}{}}
\newlabel{categories-section-localization@cref}{{[originalsection][0][]04VB}{[1][0][]source}{}{}{}}
\makeatother
\title{Classification of Yoneda extensions by derived Ext without enough projectives or injectives}
\author{The Stacks Project Authors}
\date{stacks-project:derived.tex:local-snapshot-20261011; GNU FDL 1.2 or later}
\begin{document}
\maketitle
\noindent\textit{Editorial note.} No assumption of enough injectives or projectives is made. The original common-extension construction, including the concluding elementary commutative-diagram verification, is retained.\par\medskip
\noindent\textit{Editorial scope.} One original target with its local construction proved in full. Classical category primitives are explicitly allowed.\par\medskip
\noindent\textit{Allowed standard prerequisites (editorial).}\begin{itemize}
\item Derived-category calculus of fractions. Morphisms in D(A) are roofs in the homotopy category with quasi-isomorphism denominators; equality admits a common denominator after homotopy adjustments. Standard good truncations and pushouts in an abelian category may be used.
\end{itemize}\medskip
\ifdefined\hypertarget\hypertarget{sourcepointer-1}{}\fi
\noindent\textit{Source reference.} Original section 04VB: focused-seeds/ACL/sources/stacks-snapshot-20261011/categories.tex, source lines 3893--3899. Standard localization roofs, not the target extension classification. Stable Stacks tag displayed.\par\medskip
\begin{definition}
\label{definition-yoneda-extension}
Let $\mathcal{A}$ be an abelian category. Let $A, B \in \Ob(\mathcal{A})$.
For $i \geq 1$ a degree $i$ {\it Yoneda extension} of $B$ by $A$ is an
exact sequence
$$
E : 0 \to A \to Z_{i - 1} \to Z_{i - 2} \to \ldots \to Z_0 \to B \to 0
$$
in $\mathcal{A}$. We say two Yoneda extensions $E$ and $E'$ of the same degree
are {\it equivalent} if there exists a commutative diagram
$$
\xymatrix{
0 \ar[r] & A \ar[r] & Z_{i - 1} \ar[r] & \ldots \ar[r] &
Z_0 \ar[r] & B \ar[r] & 0 \\
0 \ar[r] &
A \ar[r] \ar[u]^{\text{id}} \ar[d]_{\text{id}} &
Z''_{i - 1} \ar[r] \ar[u] \ar[d] &
\ldots \ar[r] &
Z''_0 \ar[r] \ar[u] \ar[d] &
B \ar[r] \ar[u]_{\text{id}} \ar[d]^{\text{id}} & 0 \\
0 \ar[r] & A \ar[r] & Z'_{i - 1} \ar[r] & \ldots \ar[r] &
Z'_0 \ar[r] & B \ar[r] & 0
}
$$
where the middle row is a Yoneda extension as well.
\end{definition}

\medskip\noindent
Let $\mathcal{A}$ be an abelian category with objects $A$, $B$.
Given a Yoneda extension
$E : 0 \to A \to Z_{i - 1} \to Z_{i - 2} \to \ldots \to Z_0 \to B \to 0$
we define an associated element $\delta(E) \in \Ext^i_\mathcal{A}(B, A)$
as the morphism $\delta(E) = fs^{-1} : B[0] \to A[i]$ where
$s$ is the quasi-isomorphism
$$
(\ldots \to 0 \to A \to Z_{i - 1} \to \ldots \to Z_0 \to 0 \to \ldots)
\longrightarrow
B[0]
$$
and $f$ is the morphism of complexes
$$
(\ldots \to 0 \to A \to Z_{i - 1} \to \ldots \to Z_0 \to 0 \to \ldots)
\longrightarrow
A[i]
$$
We call $\delta(E) = fs^{-1}$ the {\it class} of the Yoneda extension.
It turns out that this class characterizes the equivalence class
of the Yoneda extension.

\begin{lemma}
\label{lemma-yoneda-extension}
Let $\mathcal{A}$ be an abelian category with objects $A$, $B$ and
let $i \geq 1$.
Any element in $\Ext^i_\mathcal{A}(B, A)$ is $\delta(E)$
for some degree $i$ Yoneda extension of $B$ by $A$.
Given two Yoneda extensions $E$, $E'$ of the same degree
then $E$ is equivalent to $E'$ if and only if $\delta(E) = \delta(E')$.
\end{lemma}

\begin{proof}
Let $\xi : B[0] \to A[i]$ be an element of $\Ext^i_\mathcal{A}(B, A)$.
We may write $\xi = f s^{-1}$ for some quasi-isomorphism
$s : L^\bullet \to B[0]$ and map $f : L^\bullet \to A[i]$.
After replacing $L^\bullet$ by $\tau_{\leq 0}L^\bullet$ we may assume
that $L^j = 0$ for $j > 0$. Picture
$$
\xymatrix{
L^{- i - 1} \ar[r] & L^{-i} \ar[r] \ar[d] & \ldots \ar[r] &
L^0 \ar[r] & B \ar[r] & 0 \\
& A
}
$$
Then setting $Z_{i - 1} = (L^{- i + 1} \oplus A)/L^{-i}$ and
$Z_j = L^{-j}$ for $j = i - 2, \ldots, 0$ we see that we obtain a
degree $i$ extension $E$ of $B$ by $A$ whose class $\delta(E)$ equals
$\xi$.

\medskip\noindent
It is immediate from the definitions that equivalent Yoneda extensions
have the same class. Suppose that
$E : 0 \to A \to Z_{i - 1} \to Z_{i - 2} \to \ldots \to Z_0 \to B \to 0$ and
$E' : 0 \to A \to Z'_{i - 1} \to Z'_{i - 2} \to \ldots \to Z'_0 \to B \to 0$
are Yoneda extensions with the same class.
By construction of $D(\mathcal{A})$ as the localization
of $K(\mathcal{A})$ at the set of quasi-isomorphisms, this means there
exists a complex $L^\bullet$ and quasi-isomorphisms
$$
t : L^\bullet \to
(\ldots \to 0 \to A \to Z_{i - 1} \to \ldots \to Z_0 \to 0 \to \ldots)
$$
and
$$
t' : L^\bullet \to
(\ldots \to 0 \to A \to Z'_{i - 1} \to \ldots \to Z'_0 \to 0 \to \ldots)
$$
such that $s \circ t = s' \circ t'$ and $f \circ t = f' \circ t'$, see
Categories, Section \ref{categories-section-localization}.
Let $E''$ be the degree $i$ extension of $B$ by $A$ constructed from
the pair $L^\bullet \to B[0]$ and $L^\bullet \to A[i]$ in the first
paragraph of the proof. Then the reader sees readily that there exists
``morphisms'' of degree $i$ Yoneda extensions $E'' \to E$ and $E'' \to E'$
as in the definition of equivalent Yoneda extensions (details omitted).
This finishes the proof.
\end{proof}

\end{document}
